\documentclass[10pt,a4paper]{amsart}
\usepackage[margin=19mm]{geometry}
\usepackage{amsmath,amssymb,mathtools}
\usepackage[scr=rsfs,cal=euler]{mathalfa}
\usepackage{xfrac}
\usepackage[unicode,hidelinks,bookmarks=false]{hyperref}
\newcommand{\ie}{i.e.\ }
\newcommand{\cf}{cf.\ }
\newcommand{\resp}{respectively\ }
\newcommand{\Subsection}{\subsection}
\providecommand{\nobreakdash}{\nobreak\hbox{-}}
\setlength{\emergencystretch}{2em}
\multlinegap0pt
\usepackage{bm}
\usepackage{bbm}
\usepackage{dsfont}
\usepackage{xspace}
\usepackage{cancel}
\usepackage{mathtools}
\usepackage{amsthm}
\makeatletter
\@ifundefined{thm}{\newtheorem{thm}{Theorem}}{}
\@ifundefined{lem}{\newtheorem{lem}[thm]{Lemma}}{}
\@ifundefined{prop}{\newtheorem{prop}[thm]{Proposition}}{}
\makeatother
\DeclareMathOperator{\Ad}{Ad}
\DeclareMathOperator{\Irr}{Irr}
\DeclareMathOperator{\Rep}{Rep}
\DeclareMathOperator{\SL}{SL}
\DeclareMathOperator{\Sh}{Sh} 
\newcommand{\mbf}{\mathbf}
\newcommand{\mb}{\mathbb}
\newcommand{\mc}{\mathcal}
\newcommand{\mf}{\mathfrak}
\newcommand{\mr}{\mathrm}
\newcommand{\ol}{\overline}
\newcommand{\wh}{\widehat}
\newcommand{\wt}{\widetilde}
\title[Finite Langlands correspondence]{Finite Langlands correspondence}
\author{Naoki Imai}
\date{Source: arXiv:2508.15101v1 (CC BY 4.0)}
\begin{document}
\maketitle

\noindent\emph{Source and licence.} Finite Langlands correspondence. Version arXiv:2508.15101v1 (CC BY 4.0), distributed under the Creative Commons Attribution 4.0 International licence, as recorded in the arXiv licence metadata for this exact version and shown on the abstract page https://arxiv.org/abs/2508.15101v1. Full rights notice: licenses/arxiv-2508.15101-CC-BY-4.0.txt.\par\smallskip

\noindent\textit{Editorial scope.} Counted unit: the Thm whose source label is \texttt{thm:LCfin} in the article (source0.tex lines 1064--1073, source label \texttt{thm:LCfin}), delivered together with the article's own complete original proof of it (source0.tex lines 1074--1130, one \texttt{proof} environment of 57 lines). No mathematics is written, repaired or re-derived: statement and proof are the article's own text line for line. Required context retained uncounted: lem:ccexist (lines 494--499); thm:RepGcateq (lines 548--559); lem:ZquotOmega (lines 687--693); prop:bijSLWD (lines 1014--1032). Every definition, display and label of those lines is retained unchanged. Omitted from this package and not redistributed: 1--493, 500--547, 560--686, 694--1013, 1033--1063, 1131--1429. Nothing inside a retained excerpt is omitted or altered: every retained body is delivered exactly as the article's lines are written, so no omission receipt of any kind accompanies this package. Citations: the retained excerpts carry no citation command at all, so no bibliography entry of the article is transcribed and no reference is redistributed. Reference adaptation: none was needed and none was made; every reference inside the delivered unit resolves inside it, so no unresolved reference or citation is produced. Printed slips kept literal: no printed word, symbol, hypothesis or number of the counted statement or of its proof was altered. Layout and class: the article's own class is \texttt{article} with packages including latexsym, amsmath, amssymb, amsthm, amscd, mathrsfs, xy, graphicx, comment, arydshln, mathtools; the wrapper is the lane's pinned \texttt{amsart} editorial wrapper at 10pt on A4, which restates the theorem-like environments the retained text needs and adds the article's own macro definitions that the retained text needs (\texttt{\textbackslash Ad}, \texttt{\textbackslash Irr}, \texttt{\textbackslash Rep}, \texttt{\textbackslash SL}, \texttt{\textbackslash Sh}, \texttt{\textbackslash mb}, \texttt{\textbackslash mbf}, \texttt{\textbackslash mc}, \texttt{\textbackslash mf}, \texttt{\textbackslash mr}, \texttt{\textbackslash ol}, \texttt{\textbackslash wh}, \texttt{\textbackslash wt}), with extra wrapper packages (bm, bbm, dsfont, xspace, cancel, mathtools). The delivered numbering is the unit's own. Assets: no retained span contains a figure environment, a graphics-inclusion command, a PDF-page inclusion, a table or a TikZ picture; no image file is copied into this package, the package declares no asset dependency and no image file is redistributed with it. Licence: version arXiv:2508.15101v1 of this article is distributed under the Creative Commons Attribution 4.0 International licence, as recorded in the arXiv licence metadata for this exact version and shown on the abstract page https://arxiv.org/abs/2508.15101v1. A full rights notice is delivered at licenses/arxiv-2508.15101-CC-BY-4.0.txt. Characters: every retained excerpt is pure ASCII, so the delivered engine, XeLaTeX with its default Latin Modern text font, reports no missing character.
\begin{lem}\label{lem:ccexist}
	Let $\mf{c}$ be a unipotent parameter of $G$ with respect to a semisimple parameter $\mf{o}$ of $G$. 
	Let $\mc{L}$ be a representative of $\mf{o}$. 
	Then there is a two-sided cell 
	$\mbf{c} \subset (W_{\mc{L}})^{\circ}$ such that $\mf{c} \cap ((W_{\mc{L}})^{\circ} \times \{ \mc{L} \})=\Omega_{\mc{L}}  \mbf{c} \times \{\mc{L}\}$. 
\end{lem}

\begin{thm}\label{thm:RepGcateq}
	We fix a Whittaker datum of $G^{\circ}$. 
	Let $\mf{B}_{\mbf{c},\Omega_{\mbf{c}}}^{\circ}$ be a set of representative of 
	$\mf{B}_{\mbf{c}}^{\circ}/\Ad_{\sigma}(\Omega_{\mbf{c}})$. 
	Then there is a natural equivalence 
	\[
	\Rep (G(k)) \cong \bigoplus_{\mf{o}} \bigoplus_{\mf{c}}
	\bigoplus_{\beta \in \mf{B}_{\mbf{c},\Omega_{\mbf{c}}}^{\circ}} \Sh^{\mc{G}_{\mbf{c}},\tau_{\beta}}(\mc{G}_{\mbf{c}})^{\Omega_{\mbf{c},\beta}}
	\]
	of abelian categories, where $\mf{o}$ runs thorough the semisimple parameters of $G$, $\mf{c}$ runs thorough the unipotent parameters of $G$ with respect to $\mf{o}$ and 
	we take $\mbf{c}$ for each $\mf{c}$ as Lemma \ref{lem:ccexist}. 
\end{thm}

\begin{lem}\label{lem:ZquotOmega}
We have an equivalence  
\begin{equation}\label{eq:ZOmega}
	\Rep (Z_{\ol{A}_{\mc{H}}(u_{\mbf{c}})}(g \ol{h_{\beta} \dot{w}^{\beta} \sigma_q})) \cong \Rep (Z_{\mc{G}_{\mbf{c}}}(g \tau_{\beta}))^{\Omega_{\mbf{c},\beta}} 
\end{equation}
of abelian categories. 
\end{lem}

\begin{prop}\label{prop:bijSLWD}
Sending an L-parameter 
$\psi$ of $\SL_2$-type to $\varphi$ defined by 
\begin{equation}
 \varphi (a,w)=\psi \left( 
\begin{pmatrix}
1 & a \\ 0 & 1 
\end{pmatrix}
\begin{pmatrix}
\lvert w \rvert^{1/2} & 0 \\ 0 & \lvert w \rvert^{-1/2} 
\end{pmatrix}, w \right) , 
\end{equation}
we have bijections 
$\Psi_{\Lambda}(G) \cong \Phi_{\Lambda}(G)$ 
and 
$\Psi_{\Lambda}(G)_{\mr{sp}} \cong \Phi_{\Lambda}(G)_{\mr{sp}}$. 
Further we have a natural isomorphism 
$A_{\psi} \cong A_{\varphi}$ for $\psi$ and $\varphi$ as above. 
\end{prop}

\begin{thm}\label{thm:LCfin}
We fix a Whittaker datum of $G$. 
Then we have a natural map 
	\[
	\mc{L}_G \colon \Irr_{\ol{\mb{Q}}_{\ell}}(G(k)) \to 
	\Phi_{\ol{\mb{Q}}_{\ell}}(G)_{\mr{sp}},  
	\]
and a natural bijection between $\mc{L}_G^{-1}(\varphi)$ and 
$\Irr_{\ol{\mb{Q}}_{\ell}}(A_{\varphi})$ for $\varphi \in \Phi_{\ol{\mb{Q}}_{\ell}}(G)_{\mr{sp}}$. 
\end{thm}

\begin{proof}
We construct a map 
	\begin{equation}
	\mc{L}_G^{\Psi} \colon \Irr_{\ol{\mb{Q}}_{\ell}}(G(k))  \to 
		\Psi_{\ol{\mb{Q}}_{\ell}}(G)_{\mr{sp}}. 
	\end{equation}
Let $\pi \in \Irr_{\ol{\mb{Q}}_{\ell}}(G(k))$. 
Let $\mc{F}_{\pi}$ be an irreducible object of 
$\Sh^{\mc{G}_{\mbf{c}},\tau_{\beta}}(\mc{G}_{\mbf{c}})^{\Omega_{\mbf{c},\beta}}$ corresponding to $\pi$ under 
Theorem \ref{thm:RepGcateq} 
for some $\mf{o}$, $\mf{c}$ and 
$\beta \in \mf{B}_{\mbf{c},\Omega_{\mbf{c}}}^{\circ}$. 
This gives an element $g \in \mc{G}_{\mbf{c}}$ and a representaion of $Z_{\mc{G}_{\mbf{c}}}(g\sigma_{\beta})$ with $\Omega_{\mbf{c},\beta}$-equivariant structure. 

Let $\mc{L}$ be a representative of $\mf{o}$. 
Recall that $\mc{H}=Z_{\wh{G}}(\phi_{\mc{L},0})$. 
Let $u_{\mbf{c}} \in \mc{H}^{\circ}$ an element of the unipotent conjugacy classs corresponding to $\mbf{c}$. 
We have an identification 
$\ol{A}_{\mc{H}^{\circ}}(u_{\mbf{c}}) \cong \mc{G}_{\mbf{c}}$. 
Let $\wt{g} \in Z_{\mc{H}^{\circ}}(u_{\mbf{c}})$ be a lift of $g$. 

We define $\psi$ by 
\[	 
 \psi|_{I_k}=\phi_{\mc{L},0}, \quad 
\psi \begin{pmatrix}
	1 & 1 \\ 0 & 1 
\end{pmatrix}=u_{\mbf{c}}, \quad  
\psi (\sigma_q) = \wt{g} h_{\beta} \dot{w}^{\beta} \rtimes \sigma_q. 
\]
	We note that 
\begin{align*}
	\Ad(\psi (\sigma_q))(\phi_{\mc{L},0})=\phi_{\mc{L},0}^q. 
\end{align*}
We also have 
$\Ad (\psi (\sigma_q))(u_{\mbf{c}})=u_{\mbf{c}}$. 
We put $\mc{L}_G^{\Psi} (\pi)=\psi$. 

We have natural bijections 
\begin{equation}\label{eq:Apicent}
	A_{\psi}\cong Z_{\ol{A}_{\mc{H}}(u_{\mbf{c}})}(\ol{\psi(\sigma_q)}) \cong Z_{\ol{A}_{\mc{H}}(u_{\mbf{c}})}(g \ol{h_{\beta} \dot{w}^{\beta} \sigma_q}).
\end{equation}
By \eqref{eq:Apicent} and Lemma \ref{lem:ZquotOmega}, 
we have 
\begin{equation}\label{eq:ApsiZOmega}
\Rep (A_{\psi}) \cong \Rep (Z_{\mc{G}_{\mbf{c}}}(g \tau_{\beta}))^{\Omega_{\mbf{c},\beta}} . 
\end{equation}
By the construction and \eqref{eq:ApsiZOmega}, the fiber $(\mc{L}_G^{\Psi})^{-1}(\psi)$ 
is parametrized by 
the isomorphism classes of the irreducible representations of $A_{\psi}$. 

We define $\mc{L}_G$ 
as the composite of $\mc{L}_G^{\Psi}$ 
and 
the natural bijection 
$\Psi_{\ol{\mb{Q}}_{\ell}}(G)_{\mr{sp}} \cong 
\Phi_{\ol{\mb{Q}}_{\ell}}(G)_{\mr{sp}}$ in Proposition \ref{prop:bijSLWD}. 
\end{proof}

\end{document}
